% answers_short.tex —— 2003 年·春季北京卷·文科 · 简短答案速查
% 由 tools/make_short_answers.py 自 answers.tex 的 \gksolhead/\gkanswer 提取，
% 并与 2003/2003_数学真题_春季北京文_答案要点.md 逐题比对后生成，勿手改（改源 MD 后重跑）。

\begin{gknotice}
\textbf{使用说明}：本节只给最终答案便于快速核对，完整推导见解析版。答案为编写者独立推导，\textbf{非官方标准答案}。
\end{gknotice}

\gkpart{单选题}
\begin{center}\begin{tabular}{ccccccccccccc}
\toprule
题号 & 1 & 2 & 3 & 4 & 5 & 6 & 7 & 8 & 9 & 10 & 11 & 12 \\
\midrule
答案 & A & D & D & C & C & A & B & C & D & D & B & B \\
\bottomrule
\end{tabular}\end{center}

\gkpart{填空题}
\begin{itemize}[leftmargin=2.6em]
  \item[\textbf{13.}] $\pi$
  \item[\textbf{14.}] $\dfrac{2}{\sqrt3}$
  \item[\textbf{15.}] $140,\ 85$
  \item[\textbf{16.}] $2\sqrt3$
\end{itemize}

\gkpart{解答题}
\begin{itemize}[leftmargin=2.6em]
  \item[\textbf{17.}] $(3,+\infty)$
  \item[\textbf{18.}] 定义域为 $\{x\in\mathbb R\mid \cos2x\ne0\}=\mathbb R\setminus\{\dfrac{\pi}{4}+\dfrac{k\pi}{2}:k\in\mathbb Z\}$；$f$ 为偶函数；值域为 $[-1,2]\setminus\{\dfrac12\}$。
  \item[\textbf{19.}] （1）$\dfrac{2}{3}$；（2）证明见解析；（3）$\dfrac{\sqrt5}{2}$。
  \item[\textbf{20.}] 当 $a=1$ 时轨迹为直线 $x=0$；当 $a\ne1$ 时轨迹为圆，圆心 $\left(\dfrac{c(1+a^2)}{a^2-1},0\right)$，半径 $\dfrac{2ac}{|a^2-1|}$。
  \item[\textbf{21.}] （1）$88$ 辆；（2）月租金定为 $4050$ 元时，月收益最大，为 $307050$ 元。
  \item[\textbf{22.}] （1）$\dfrac{a_{n+1}}{a_n}=\dfrac19$；（2）$\dfrac{3\pi l^2}{32}$。
\end{itemize}
